% SOSP 2026 Paper Template
% The 32nd ACM Symposium on Operating Systems Principles
% Conference: September 30, 2026
%
% Format: ACM SIGPLAN, ≤12 pages technical content (excluding references)
%         A4 or US letter, 178×229mm (7×9") text block
%         Two-column, 8mm separation, 10pt on 12pt leading
%
% Official CFP: https://sigops.org/s/conferences/sosp/2026/cfp.html
% ACM Template: https://www.acm.org/publications/proceedings-template
%
% IMPORTANT NOTES:
% - Double-blind review (use paper ID, not author names)
% - Anonymized system/project name required (different from arXiv/talks)
% - Optional Artifact Evaluation after acceptance
% - Author response period available
%   → Limited to: correcting factual errors + addressing questions
%   → NO new experiments or additional work
%   → Keep under 500 words
% - Supplementary material allowed (reviewers not required to read)
% - Figures/tables readable without magnification, color encouraged
% - Pages numbered, references hyperlinked

\documentclass[sigplan,10pt]{acmart}

% Remove copyright/permission footer for submission
\renewcommand\footnotetextcopyrightpermission[1]{}
\settopmatter{printfolios=true}

% Remove ACM reference format for submission
\setcopyright{none}
\renewcommand\acmConference[4]{}
\acmDOI{}
\acmISBN{}

% Recommended packages
\usepackage{booktabs}
\usepackage{xspace}
\usepackage{subcaption}

% Custom commands
\newcommand{\system}{SystemName\xspace}  % Replace with your anonymized system name
\newcommand{\eg}{e.g.,\xspace}
\newcommand{\ie}{i.e.,\xspace}
\newcommand{\etal}{\textit{et al.}\xspace}

\begin{document}

\title{Your Paper Title Here}

% Anonymized for submission - use paper ID
\author{Paper \#XXX}
\affiliation{}

% Camera-ready:
% \author{Author One}
% \affiliation{%
%   \institution{University/Company}
%   \city{City}
%   \country{Country}}
% \email{email@example.com}

\begin{abstract}
% Write your abstract here.
Your abstract goes here. Summarize your contribution to computer systems
software design, implementation, analysis, evaluation, or deployment.
\end{abstract}

\maketitle
\pagestyle{plain}

\section{Introduction}
\label{sec:intro}

% SOSP encourages groundbreaking work in significant new directions.
% Evaluation criteria for papers addressing new problems may differ
% from those in established areas.
%
% A good SOSP paper:
% - Motivates a significant problem
% - Proposes and implements an interesting, compelling solution
% - Demonstrates the practicality and benefits of the solution
% - Draws appropriate conclusions
% - Clearly describes contributions and advances beyond previous work

Your introduction goes here. Clearly state the problem, your approach,
and your contributions.

\section{Background and Motivation}
\label{sec:background}

Provide necessary background and motivation for your systems work.

\section{Design}
\label{sec:design}

Describe the design of your system.

\section{Implementation}
\label{sec:implementation}

Describe implementation details.

\section{Evaluation}
\label{sec:evaluation}

Present your evaluation results with real workloads.

\subsection{Experimental Setup}

Describe your experimental setup, workloads, and baselines.

\subsection{End-to-End Performance}

Present end-to-end system performance results.

\subsection{Microbenchmarks}

Present targeted microbenchmarks to explain specific design decisions.

\section{Discussion}
\label{sec:discussion}

Discuss limitations, lessons learned, and future directions.

\section{Related Work}
\label{sec:related}

Discuss related work in operating systems and systems software.

\section{Conclusion}
\label{sec:conclusion}

Summarize your contributions and key findings.

% Acknowledgments (only in camera-ready, remove for submission)
% \begin{acks}
% We thank...
% \end{acks}

\bibliographystyle{ACM-Reference-Format}
\bibliography{references}

\end{document}
